\section{Geometry of the dyadic distribution protocol}
\label{sec:peps_dyadic_geometry}
This section treats the geometric part of~\cite[Section~7]{OpenAI2026PolynomialPEPS},
September 24, 2026, \texttt{06-geometry.tex}: the bands along the edges of a
repainted dyadic square, the angular separation of~\cite[Lemma~7.2]{OpenAI2026PolynomialPEPS}
(lines 254--319), the point-treatment estimates (lines 357--415), and the counts
of the padded dyadic hierarchy behind items 1--3
of~\cite[Proposition~7.1]{OpenAI2026PolynomialPEPS} (lines 12--53 and 633--652).
The protocol of Proposition~7.1 itself, whose changes are implemented by the
information estimate and the contraction gates of the earlier sections, is not
treated here. Distances are sup distances in the plane, and labels are elements
of an arbitrary set; the labels $A$, $B$, $C$ below are permitted to coincide.

\begin{definition}[Edge coordinates, bands and the lens]
    \label{def:peps_dyadic_edge_band}
    \lean{TNLean.PEPS.Approximation.SquareEdge, TNLean.PEPS.Approximation.bandWidth, TNLean.PEPS.Approximation.edgeBand, TNLean.PEPS.Approximation.closedEdgeBand, TNLean.PEPS.Approximation.bandCoord, TNLean.PEPS.Approximation.edgeLens, TNLean.PEPS.Approximation.edgeNeighbour, TNLean.PEPS.Approximation.edgeMarkDist, TNLean.PEPS.Approximation.cornerMarkDist}
    \leanok
    Let $S=[0,n]^2$. For an edge of $S$ let $s$ be the coordinate parallel to the
    edge, measured from one endpoint, and $d$ the normal coordinate, positive
    into $S$. Put
    \begin{align}
        w(s)=10^{-3}\min(s,n-s),\quad x=d/w(s).
        \label{eq:pdg_band}
    \end{align}
    The band $\{\alpha<x<\beta\}$ is the set of points with $0<s<n$ and
    $\alpha w(s)<d<\beta w(s)$; its closed version has $0\le s\le n$ and
    $\alpha w(s)\le d\le\beta w(s)$. The lens of the edge is the band
    $Y=\{-7<x<-7/2\}$, and the exterior neighbouring square is
    $\{0<s<n,\ -n<d<0\}$. For a point $y$, $d_{\mathcal V}(y)$ is its distance from
    the two endpoints of the edge, or from the four corners of $S$.
\end{definition}

\begin{lemma}[Disjointness of the edge bands]
    \label{lem:peps_dyadic_band_disjoint}
    \lean{TNLean.PEPS.Approximation.disjoint_edgeBand, TNLean.PEPS.Approximation.edgeBand_neg_subset_edgeNeighbour, TNLean.PEPS.Approximation.dist_le_of_mem_closedEdgeBand}
    \leanok
    \uses{def:peps_dyadic_edge_band}
    The bands $\{-8<x<2\}$ along distinct edges of $S$ are disjoint, and the band
    $\{-8<x<0\}$ lies in the exterior neighbouring square of its edge. Two points
    of a closed band with $-8\le\alpha\le\beta\le8$ are at distance at most $n$.
\end{lemma}

\begin{proof}\leanok
    Write $w\le s/1000$ and $w\le(n-s)/1000$. For opposite edges the bands lie
    within $n/250$ of their edges. For adjacent edges a common point at distances
    $r_1,r_2>0$ from the two edge lines would satisfy $r_1<0.002\,r_2$ and
    $r_2<0.002\,r_1$. On $\{-8<x<0\}$ one has $0>d>-8w\ge-n/250$. The diameter
    bound follows from $|d|\le8w\le n/250$ in each closed band.
\end{proof}

\begin{definition}[Normal words of the edge construction]
    \label{def:peps_dyadic_normal_words}
    \lean{TNLean.PEPS.Approximation.bandWord, TNLean.PEPS.Approximation.edgeStartWord, TNLean.PEPS.Approximation.auxWord, TNLean.PEPS.Approximation.mainWordOne, TNLean.PEPS.Approximation.edgeOperations, TNLean.PEPS.Approximation.BandOperation}
    \leanok
    \uses{def:peps_dyadic_edge_band}
    A normal word lists labels in increasing $x$ separated by interfaces; at an
    interface the label on its right is taken. An edge operation starts from the
    main word $C\mid A\mid B$ with interfaces $0,1$. On an auxiliary sheet, uniform
    at $C$, the births of $B$ in $(-6,-1)$, of $A$ in $(-5,-2)$ and of $C$ in
    $(-4,-3)$ give
    \begin{align}
        C\mid B\mid A\mid C\mid A\mid B\mid C\quad\text{at } -6,-5,-4,-3,-2,-1.
        \label{eq:pdg_aux_word}
    \end{align}
    After the lens exchange the main word is $C\mid B\mid A\mid C\mid A\mid B$ at
    $-6,-5,-4,0,1$. The $C$ band $(-4,0)$ then dies into $A$, the $A$ band
    $(-5,1)$ dies into $B$, $C$ is born in $(-3,0)$, and the $B$ band $(-6,-3)$
    dies into $C$, leaving $C\mid B$ at $0$. Each of these seven births and deaths
    changes its word only on a closed interval, and its surrounding word (the word
    before a birth, after a death) carries one label on a strictly larger open
    interval inside $(-8,2)$, with margins at least $1/2$.
\end{definition}

\begin{lemma}[The lens exchange produces the main word]
    \label{lem:peps_dyadic_lens_word}
    \lean{TNLean.PEPS.Approximation.lens_exchange_word}
    \leanok
    \uses{def:peps_dyadic_normal_words}
    Replacing the starting main word by the auxiliary word (\ref{eq:pdg_aux_word})
    on $-7<x<-7/2$ gives the main word $C\mid B\mid A\mid C\mid A\mid B$ at
    $-6,-5,-4,0,1$, at every value of $x$.
\end{lemma}

\begin{proof}\leanok
    Both thresholds $-7$ and $-7/2$ lie in $C$ intervals of both words; compare
    the two words on each interval.
\end{proof}

\begin{lemma}[Conical band estimate]
    \label{lem:peps_dyadic_conical_band}
    \lean{TNLean.PEPS.Approximation.mem_edgeBand_of_dist_lt, TNLean.PEPS.Approximation.mem_edgeBand_of_dist_lt_markDist, TNLean.PEPS.Approximation.edgeMarkDist_le_of_mem_closedEdgeBand}
    \leanok
    \uses{def:peps_dyadic_edge_band}
    Let $a_0=1/5000$, $\alpha\le\beta$, $\alpha'+1/2\le\alpha$, $\beta+1/2\le\beta'$
    and $-8\le\alpha'$, $\beta'\le8$. If $y$ lies in the closed band
    $\{\alpha\le x\le\beta\}$ and $|y-z|<a_0\min(n,d_{\mathcal V}(y))$, then $z$
    lies in the open band $\{\alpha'<x<\beta'\}$.
\end{lemma}

\begin{proof}\leanok
    Put $m=\min(s,n-s)$ at $y$, so $w=m/1000$ and $|d|\le m$; hence
    $d_{\mathcal V}(y)\le m$. If $|y-z|<m/5000$, the parallel and normal
    coordinates move by less than $m/5000$, and $w$ is $10^{-3}$-Lipschitz, so
    $|\beta'(w(s_z)-w(s_y))|\le8m/(5\cdot10^{6})$. Then
    $d_z<\beta w(s_y)+m/5000\le\beta' w(s_z)-m/2000+m/5000+8m/(5\cdot10^6)$, and
    symmetrically at $\alpha'$.
\end{proof}

\begin{theorem}[Angular separation of the central birth]
    \label{thm:peps_dyadic_central_clearance}
    \lean{TNLean.PEPS.Approximation.centralRegion, TNLean.PEPS.Approximation.centralBirth_clearance}
    \leanok
    \uses{def:peps_dyadic_edge_band}
    Let $R$ be the set of points of $(0,n)^2$ with $x>1$ for every edge, and let
    $f$ be a labelling of the plane equal to $A$ on $(0,n)^2$. For every $y$ in
    the closure of $R$ and every $z$ with $f(z)\ne A$,
    \begin{align}
        |y-z|\ge a_0\min\{n,d_{\mathcal V}(y)\},
        \label{eq:pdg_central}
    \end{align}
    where $\mathcal V$ is the set of corners of $S$ and $a_0=1/5000$. This is the
    central-birth case of~\cite[Lemma~7.2]{OpenAI2026PolynomialPEPS}, whose
    surrounding guide is $A$ on the interior of $S$.
\end{theorem}

\begin{proof}\leanok
    With $u=\min(y_1,n-y_1)$ and $v=\min(y_2,n-y_2)$, the four conditions $d\ge w$
    give $v\ge u/1000$ and $u\ge v/1000$. The nearest corner is at distance
    $\max(u,v)\le1000\min(u,v)$, so a point within $a_0\min(n,d_{\mathcal V}(y))$
    of $y$ is within $\min(u,v)/5$ of $y$ and lies in $(0,n)^2$.
\end{proof}

\begin{theorem}[Angular separation of the edge births and deaths]
    \label{thm:peps_dyadic_band_clearance}
    \lean{TNLean.PEPS.Approximation.BandOperation.clearance, TNLean.PEPS.Approximation.auxBirthB, TNLean.PEPS.Approximation.auxBirthA, TNLean.PEPS.Approximation.auxBirthC, TNLean.PEPS.Approximation.mainDeathC, TNLean.PEPS.Approximation.mainDeathA, TNLean.PEPS.Approximation.mainBirthC, TNLean.PEPS.Approximation.mainDeathB}
    \leanok
    \uses{lem:peps_dyadic_conical_band, def:peps_dyadic_normal_words}
    Consider one of the seven elementary births and deaths of an edge operation,
    and labellings $f_b$, $f_a$ and $f_s$ of the plane whose normal words on
    $\{-8<x<2\}$ are its words before, after and surrounding, with $f_b=f_a$ off
    that band. Let $P_\circ$ be its surrounding label. For every $y$ in the closure
    of $\{f_b\ne f_a\}$ and every $z$ with $f_s(z)\ne P_\circ$,
    $|y-z|\ge a_0\min\{n,d_{\mathcal V}(y)\}$, with $\mathcal V$ the two endpoints
    of the edge. These are the edge-band cases of~\cite[Lemma~7.2]{OpenAI2026PolynomialPEPS}.
\end{theorem}

\begin{proof}\leanok
    The set $\{f_b\ne f_a\}$ lies in the closed band of the changed interval,
    which is closed. If $|y-z|<a_0\min\{n,d_{\mathcal V}(y)\}$, the conical band
    estimate puts $z$ in the open band of the surrounding sector, where $f_s$
    equals $P_\circ$.
\end{proof}

\begin{theorem}[Angular separation of the lens exchange]
    \label{thm:peps_dyadic_lens_clearance}
    \lean{TNLean.PEPS.Approximation.lensExchange_clearance}
    \leanok
    \uses{lem:peps_dyadic_conical_band, def:peps_dyadic_normal_words}
    Let the main labelling be $C$ on the exterior neighbouring square, and let the
    auxiliary labelling have the word (\ref{eq:pdg_aux_word}) on $\{-8<x<2\}$ and
    be $C$ off it. Every point $y$ in the closure of the positions that are not
    $C$ on both labellings satisfies
    $|y-z|\ge a_0\min\{n,d_{\mathcal V}(y)\}$ for every $z$ on the boundary of the
    lens $Y$. This is the exchange case of~\cite[Lemma~7.2]{OpenAI2026PolynomialPEPS}.
\end{theorem}

\begin{proof}\leanok
    A boundary point of $Y$ is an endpoint of the edge, where the bound is
    immediate, or satisfies $d=-7w$ or $d=-\tfrac72w$ with $0<s<n$. If $y$ lies
    outside the exterior neighbouring square, then $d_{\mathcal V}(y)\le
    d_{\mathcal V}(z)+|y-z|$ and $d_{\mathcal V}(z)\le m_z=\min(s_z,n-s_z)$; a
    distance below $a_0d_{\mathcal V}(y)$ is then below $m_z/2000$, which keeps
    $y$ in that square because $d_z\le-\tfrac72m_z/1000$. If $y$ lies in one of
    the closed auxiliary bands $\{-6\le x\le-4\}$ or $\{-3\le x\le-1\}$, apply the
    conical band estimate with the open bands $\{-\tfrac{13}2<x<-\tfrac72\}$ or
    $\{-\tfrac72<x<-\tfrac12\}$, neither of which meets the two boundary curves.
\end{proof}

\begin{lemma}[Logarithmic floor and outer-hole enlargement]
    \label{lem:peps_dyadic_hole_enlargement}
    \lean{TNLean.PEPS.Approximation.clearance_floor, TNLean.PEPS.Approximation.tapered_clearance_of_center}
    \leanok
    If $a\min(n,d_{\mathcal V}(y))\le D$, $t\le d_{\mathcal V}(y)$ and $t\le n$,
    then $a\max\{t,\min(n,d_{\mathcal V}(y))\}\le D$. Let $g$ and $h$ be
    $1$-Lipschitz functions, $a,t>0$ and $2\epsilon_0(a+1)\le a/2$. If a center $c$
    satisfies $g(c)\ge aQ_c$ with $Q_x=\max\{t,\min(n,h(x))\}$, then every $x$ with
    $|x-c|\le2\epsilon_0t$ satisfies $g(x)\ge\tfrac a2Q_x$.
    This is~\cite[Section~7.4]{OpenAI2026PolynomialPEPS}, lines 357--367 and
    405--415.
\end{lemma}

\begin{proof}\leanok
    The first claim holds because $\min(n,d_{\mathcal V}(y))\ge t$. For the
    second, with $\delta=|x-c|$ one has $Q_x\le Q_c+\delta$ and
    $g(x)\ge g(c)-\delta$, so
    $g(x)\ge aQ_x-(a+1)\delta\ge aQ_x-2\epsilon_0(a+1)t\ge aQ_x-\tfrac a2Q_x$.
\end{proof}

\begin{definition}[Padded dyadic hierarchy, junctions and addresses]
    \label{def:peps_dyadic_hierarchy}
    \lean{TNLean.PEPS.Approximation.DyadicSquare, TNLean.PEPS.Approximation.DyadicJunction, TNLean.PEPS.Approximation.squareAddress}
    \leanok
    For $N=2^k$ the hierarchy consists of the aligned squares of $[0,N]^2$ of side
    $2^j$, $0\le j\le k$, indexed by $j$ and two block indices below $2^{k-j}$; the
    junctions at scale $j$ are the grid corners, indexed by two integers at most
    $2^{k-j}$. A square of side $2^j$ with block index $r$ has address coordinate
    $2^jr+2^{j-1}$ for $j\ge1$ and $r$ for $j=0$.
\end{definition}

\begin{theorem}[Counts of the hierarchy and the schedule]
    \label{thm:peps_dyadic_hierarchy_counts}
    \lean{TNLean.PEPS.Approximation.three_mul_card_dyadicSquare_add_one, TNLean.PEPS.Approximation.three_mul_card_dyadicSquare_lt, TNLean.PEPS.Approximation.three_mul_card_dyadicJunction_le, TNLean.PEPS.Approximation.three_mul_schedule_le, TNLean.PEPS.Approximation.three_mul_schedule_le_cube}
    \leanok
    \uses{def:peps_dyadic_hierarchy}
    The hierarchy has $(4^{k+1}-1)/3$ squares, and three times the number of
    junctions over all scales is at most $4^{k+2}$. If $2^k<2L$, the hierarchy has
    fewer than $16L^2/3$ squares, and a schedule using at most $M$ changes for each
    repainting of a square and each junction resize has at most $80ML^2/3$, and
    for $L\ge1$ at most $80ML^3/3$, changes.
\end{theorem}

\begin{proof}\leanok
    The scale $j$ contributes $4^{k-j}$ squares and $(2^{k-j}+1)^2\le4\cdot4^{k-j}$
    junctions, and $3\sum_{i<k+1}4^i+1=4^{k+1}$ by induction. Finally
    $4^k=(2^k)^2<4L^2$.
\end{proof}

\begin{lemma}[Addresses and locality]
    \label{lem:peps_dyadic_address_locality}
    \lean{TNLean.PEPS.Approximation.squareAddress_mem, TNLean.PEPS.Approximation.squareAddress_lt, TNLean.PEPS.Approximation.squareAddress_lt_add, TNLean.PEPS.Approximation.squareAddress_dist_lt, TNLean.PEPS.Approximation.card_blocks_meeting_le}
    \leanok
    \uses{def:peps_dyadic_hierarchy}
    An address coordinate lies in $[2^jr,2^jr+2^j)$, hence below $2^k$. If the
    blocks $[2^jr,2^j(r+1)]$ and $[2^{j'}r',2^{j'}(r'+1)]$ are within distance $g$
    and $j'\le j+1$, their address coordinates differ by less than $g+3\cdot2^j$.
    The blocks of side $b$ meeting $[V,U]$ number at most
    $\lfloor U/b\rfloor+2-\lfloor V/b\rfloor$.
\end{lemma}

\begin{proof}\leanok
    The first two claims follow from the position of the address in its block and
    $2^{j'}\le2\cdot2^j$. For the count, $r\mapsto r+1$ maps the meeting blocks
    injectively into $[\lfloor V/b\rfloor,\lfloor U/b\rfloor+1]$.
\end{proof}
